% BHU PYQ -- Manually transcribed & error-corrected
\documentclass[11pt,a4paper]{article}

\usepackage[a4paper,top=2.5cm,bottom=2.5cm,left=2.8cm,right=2.8cm,headheight=14pt]{geometry}
\usepackage{amsmath,amssymb}
\usepackage[T1]{fontenc}
\usepackage[utf8]{inputenc}
\usepackage{lmodern}
\usepackage{microtype}
\usepackage{fancyhdr}
\usepackage{enumitem}
\usepackage{hyperref}
\usepackage{lastpage}
\usepackage{parskip}

\hypersetup{colorlinks=true,urlcolor=black,linkcolor=black,
  pdftitle={CHB-101: Structure and Bonding & Organic Chemistry-I -- BHU PYQ},pdfauthor={Banaras Hindu University}}

\pagestyle{fancy}\fancyhf{}
\renewcommand{\headrulewidth}{0.4pt}\renewcommand{\footrulewidth}{0.4pt}
\fancyhead[L]{\small   \textit{Banaras Hindu University}}
\fancyhead[R]{\small   \textit{CHB-101: Structure & Bonding & Organic Chemistry-I}}
\fancyfoot[C]{\small Page  \thepage\ of \pageref{LastPage}}


\newlist{parts}{enumerate}{2}
\setlist[parts,1]{label=(\alph*),leftmargin=*,itemsep=5pt,topsep=3pt}
\setlist[parts,2]{label=(\roman*),leftmargin=*,itemsep=3pt,topsep=2pt}

\newcommand{\pts}[1]{\hfill   \textit{[#1]}}


\begin{document}


\begin{center}
  {\large\bfseries BANARAS HINDU UNIVERSITY}\\[2pt]
  {\normalsize B.Sc. (Hons.) Semester I Examination, 2022-23}\\[6pt]
  \rule{0.8\linewidth}{0.5pt}\\[6pt]
  {\large\bfseries Paper No. CHB-101: Structure \& Bonding and Organic Chemistry-I}\\[4pt]
  {\normalsize Time: 3 Hours\hfill Maximum Marks: 70}
\end{center}
\rule{\linewidth}{0.4pt}
\small



\noindent   \textit{Instructions: Answer five questions in total, including Question~1 from each section which is compulsory. Use separate answer books for Section~A and Section~B.}

\normalsize
\bigskip


\begin{center}
  {\large\bfseries SECTION A}\\[2pt]
  {\normalsize (Structure and Bonding -- Marks: 35)}
\end{center}
\rule{0.4\linewidth}{0.2pt}
\smallskip




\noindent   \textbf{Question 1.} Answer all parts of the following: \pts{5 $ \times$ 2 = 10}

\begin{parts}
  \item What kind of interaction exists among the atoms when helium is liquefied?
  \item Discuss the Electron Pool Theory in brief.
  \item What is meant by a node in an orbital?
  \item Arrange the orbitals $3 \text{s}$, $3 \text{d}$, and $4 \text{s}$ in order of increasing energy in the case of the hydrogen atom.
  \item Why is it necessary to convert the Schr\"odinger wave equation in terms of polar coordinates for a hydrogen-like atom?
\end{parts}

\medskip\rule{\linewidth}{0.2pt}




\noindent   \textbf{Question 2.}

\begin{parts}
  \item What is the expected bond order in the $ \text{Be}_2$ molecule? How would you explain the fact that beryllium exists as a metal? \pts{6}
  \item Discuss the solid-state structure of $ \text{NaCl}$ using a suitable diagram. \pts{6}
\end{parts}

\medskip\rule{\linewidth}{0.2pt}




\noindent   \textbf{Question 3.}

\begin{parts}
  \item Write the expressions for $x$, $y$, and $z$ in terms of polar coordinates for the hydrogen atom. \pts{6}
  \item Using the VSEPR model, explain the geometries and bond angles of $ \text{NH}_3$ and $ \text{NF}_3$. \pts{6}
\end{parts}

\medskip\rule{\linewidth}{0.2pt}




\noindent   \textbf{Question 4.}

\begin{parts}
  \item Derive the Schr\"odinger wave equation assuming the electron to be a standing wave and explain all the terms involved in it. \pts{6}
  \item Discuss the bonding in $ \text{H}_2^+$ molecule ion on the basis of Molecular Orbital Theory. \pts{6}
  \item How would you explain the fact that the outermost electron in potassium goes to the $4 \text{s}$ orbital rather than a $3 \text{d}$ orbital? \pts{6}
\end{parts}


\newpage

\begin{center}
  {\large\bfseries SECTION B}\\[2pt]
  {\normalsize (Organic Chemistry-I -- Marks: 35)}
\end{center}
\rule{0.4\linewidth}{0.2pt}
\smallskip




\noindent   \textbf{Question 5.} Answer all parts of the following: \pts{5 $ \times$ 2 = 10}

\begin{parts}
  \item Write appropriate reactions for the formation of   \textit{cis}-$1,2$-cyclopentanediol and   \textit{trans}-$1,2$-cyclopentanediol from cyclopentene.
  \item Draw the energy profile diagram for the following reaction:
    \[  \text{(CH}_3 \text{)}_3 \text{C-Cl} +  \text{NaOH}  o  \text{(CH}_3 \text{)}_3 \text{C-OH} +  \text{NaCl} \]
  \item Identify whether the chiral center in 2-chlorobutane is $R$ or $S$ for both enantiomers.
  \item Give the name and structure of a reagent for the selective oxidation of primary alcohols to aldehydes.
  \item Explain why the chair conformation of cyclohexane is free of angle and torsional strain.
\end{parts}

\medskip\rule{\linewidth}{0.2pt}




\noindent   \textbf{Question 6.}

\begin{parts}
  \item Discuss the base-catalysed tautomerisation of ethyl acetoacetate. \pts{6}
  \item Discuss the reactions of the following with ethyl acetoacetate, showing mechanisms: \pts{6}
  
\begin{parts}
    \item $ \text{HNO}_2$
    \item $ \text{NH}_3$
  \end{parts}
\end{parts}

\medskip\rule{\linewidth}{0.2pt}




\noindent   \textbf{Question 7.}

\begin{parts}
  \item Discuss in detail $1,2$- vs. $1,4$-addition reaction of chlorine on $1,3$-butadiene with energy profile diagrams. \pts{6}
  \item Write in detail about the esterification reaction of alcohols and compare the relative reactivity of primary, secondary, and tertiary alcohols. \pts{6}
\end{parts}

\medskip\rule{\linewidth}{0.2pt}




\noindent   \textbf{Question 8.}

\begin{parts}
  \item Compare the acidity of ethane, ethylene, and acetylene, and justify the preparation of sodium acetylide. \pts{6}
  \item Complete the following reactions and write the major product: \pts{6}
  
\begin{parts}
    \item $ \text{CH}_3 \text{-C}\equiv \text{CH} +  \text{H}_2 \text{O} \xrightarrow{ \text{dil. H}_2 \text{SO}_4,  \text{HgSO}_4} ?$
    \item $ \text{CH}_3 \text{-CH=CH}_2 +  \text{HBr} \xrightarrow{ \text{Peroxide}} ?$
  \end{parts}
\end{parts}

\vfill
\rule{\linewidth}{0.4pt}

\begin{center}\small
  \href{https://bhu-exam.vercel.app/}{Click here to visit our website}\\[4pt]
  \href{https://me-aryan.vercel.app/}{Click here to visit our contributor}\\[4pt]
  \href{https://quashan-me.vercel.app/}{Click here to visit our contributor}
\end{center}

\end{document}
